% TOP_PROBLEM: {"problemId": "problem.mahler-conjecture-on-the-volume-product-of-a-convex-body", "canonicalTitle": "Mahler Conjecture on the Volume Product of a Convex Body", "releaseRank": 302, "primaryDomain": "combinatorics_discrete_geometry", "primaryDomainLabel": "Combinatorics & discrete geometry", "displayStatus": "open", "releaseStatus": "open"}
% !TeX program = lualatex
% ATTEMPT_STATUS: unresolved
\documentclass[11pt]{article}
\usepackage[margin=25mm]{geometry}
\usepackage{fontspec}
\setmainfont{Latin Modern Roman}
\usepackage{amsmath,amssymb,mathtools}
\usepackage{unicode-math}
\setmathfont{Latin Modern Math}
\usepackage{xurl,hyperref}
\hypersetup{colorlinks=true,urlcolor=blue}
\setcounter{secnumdepth}{0}
\setlength{\parindent}{0pt}
\setlength{\parskip}{5pt}
\setlength{\emergencystretch}{3em}
\begin{document}
\section*{\texorpdfstring{Mahler Conjecture on the Volume Product of a Convex Body}{Mahler Conjecture on the Volume Product of a Convex Body}}\label{302}
\subsection{Definitions and mathematical statement}
A convex body $K\subset{\mathbb{R}}^n$ is compact, convex, and has nonempty interior. Assume $K=-K$ and define its polar by
\[
 K^\circ=\{y\in{\mathbb{R}}^n:\langle x,y\rangle\le1\text{ for every }x\in K\}.
\]
With $|\cdot|$ denoting $n$-dimensional Lebesgue volume, the symmetric Mahler conjecture asserts
\[
 |K|\,|K^\circ|\ge\frac{4^n}{n!}\qquad(n\ge1).
\]
For $K=[-1,1]^n$, its polar is $\{y:\sum_i|y_i|\le1\}$, giving equality. The product is unchanged under invertible linear changes of coordinates. The selected statement concerns origin-symmetric bodies; it is not the different conjecture about general nonsymmetric bodies and simplices, nor a classification of every equality case.
\par\smallskip\textit{Formulation and concept references:} \href{https://arxiv.org/pdf/1902.08971}{[S1]}.

\subsection{Short English statement}
Is the product of a centrally symmetric convex body's volume and its polar's volume always at least the value attained by a cube?
\subsection{Research attempt: exact polar products and the loss in a maximal-basis bound}
\paragraph{1. Scope and primary-source audit.}
The body is centrally symmetric about the origin, and the requested constant is $4^n/n!$. This is different from the nonsymmetric conjecture with its simplex constant and optimized polar center. Karasev's source [S1] addresses special hyperplane sections. The primary paper of Iriyeh and Shibata [E1], read on 27 September 2026, proves the full centrally symmetric three-dimensional inequality in Theorem 1.1. Its reference to a ``general case'' in later sections remains within dimension three. A May 2026 alternative-proof preprint [E2] has the same dimensional restriction. Neither is an all-dimensional theorem. The work below derives exact closure rules, a large equality family, and an elementary universal bound that exposes a precise factorial loss.

\paragraph{2. Linear invariance, polarity, and the one-dimensional case.}
For an invertible linear map $A$, the definition of the polar gives
\[
 (AK)^\circ=A^{-T}K^\circ.
\]
The two volumes multiply respectively by $|\det A|$ and its reciprocal, so the volume product is invariant. Also, $K\subset L$ implies $L^\circ\subset K^\circ$, and the bipolar theorem gives $(K^\circ)^\circ=K$. These facts do not make the volume product monotone under inclusion: its two factors move in opposite directions.

In dimension one, $K=[-a,a]$ and $K^\circ=[-1/a,1/a]$, so the product is four. For a box $K=\prod[-a_i,a_i]$, its polar is the weighted crosspolytope
\[
 K^\circ=\left\{y:\sum_i a_i|y_i|\leq1\right\}.
\]
Each orthant contributes a simplex of volume $1/(n!\prod_i a_i)$, giving $|K^\circ|=2^n/(n!\prod_i a_i)$. Since $|K|=2^n\prod_i a_i$, equality holds in the conjecture. Every linear image of such a box has the same equality property.

\paragraph{3. A complete product formula, including the factorial.}
Let $K\subset\mathbb R^a$ and $L\subset\mathbb R^b$ be centrally symmetric convex bodies. Their product has support function
\[
 h_{K\times L}(u,v)=h_K(u)+h_L(v).
\]
The functions $h_K,h_L$ are the gauges of $K^\circ,L^\circ$. Therefore
\[
 (K\times L)^\circ=\{(u,v):h_K(u)+h_L(v)\leq1\}.
\]
The sublevel set $\{h_K\leq r\}$ has volume $r^a|K^\circ|$. Integrating first in the gauge radius of $u$ gives
\[
 |(K\times L)^\circ|
 =a|K^\circ||L^\circ|\int_0^1r^{a-1}(1-r)^b\,dr
 =\frac{a!b!}{(a+b)!}|K^\circ||L^\circ|.
\]
The beta integral here follows by integrating by parts $b$ times, or directly from the simplex-volume formula. Multiplying by $|K\times L|=|K||L|$, one obtains
\[
 \frac{(a+b)!}{4^{a+b}}\,|K\times L||(K\times L)^\circ|
 =\left(\frac{a!}{4^a}|K||K^\circ|\right)
  \left(\frac{b!}{4^b}|L||L^\circ|\right).
\]
Thus the normalized volume product is multiplicative for Cartesian products. If the conjecture holds for both factors, it holds for their product. If both have equality, their product also has equality. Omitting the beta-integral factor would give the wrong constant already for a square built from two intervals.

\paragraph{4. Convex sums in complementary subspaces give another exact closure rule.}
Define
\[
 K\oplus_1 L=\operatorname{conv}\bigl((K\times\{0\})\cup(\{0\}\times L)\bigr).
\]
A linear functional is at most one on this convex hull exactly when it is at most one on each of the two constituent sets. Hence
\[
 (K\oplus_1L)^\circ=K^\circ\times L^\circ.
\]
Apply the preceding product formula to the polar bodies and use bipolarity. This gives the same multiplicativity of normalized volume product for $\oplus_1$. Starting with symmetric intervals and repeatedly taking products or these complementary-subspace convex sums therefore produces equality bodies in every dimension. These are the usual recursively constructed Hanner polytopes; the equality calculation itself has been proved here and does not rely on their classification.

The construction permits mixed shapes, not just cubes and crosspolytopes. For example, the product of a two-dimensional diamond and an interval has equality in dimension three. Its polar is a bipyramid over a square. In higher dimensions one may interleave the two operations in many different ways. None of this proves that every minimizer has this form, and the question does not require a classification of equality cases.

\paragraph{5. An elementary bound for every symmetric convex body.}
Choose vectors $v_1,\ldots,v_n\in K$ maximizing $D=|\det(v_1,\ldots,v_n)|$. Compactness gives a maximum and nonempty interior gives $D>0$. Write an arbitrary $x\in K$ as $x=\sum_i t_iv_i$. Replacing the $i$th column by $x$ changes the determinant to $t_i\det(v_1,\ldots,v_n)$. Maximality implies $|t_i|\leq1$ for every $i$. Central symmetry and convexity also imply that $\operatorname{conv}\{\pm v_i\}$ lies in $K$. Thus
\[
 C:=\operatorname{conv}\{\pm v_1,\ldots,\pm v_n\}
 \subset K\subset P:=\left\{\sum_i t_iv_i:|t_i|\leq1\right\}.
\]
Taking polars reverses the containments: $P^\circ\subset K^\circ\subset C^\circ$. The explicit volumes are
\[
 |C|=\frac{2^nD}{n!},\qquad |P^\circ|=\frac{2^n}{n!D}.
\]
Therefore
\[
 |K||K^\circ|\geq |C||P^\circ|=\frac{4^n}{(n!)^2}.
\]
This is a genuine universal lower bound proved from a maximal determinant. It is weaker than the requested bound by a factor $n!$. The exact loss is visible: the proof independently replaces $K$ by the small crosspolytope $C$ and $K^\circ$ by the small crosspolytope $P^\circ$. They cannot both be sharp approximations in the equality configuration when $n>1$. A sharp proof must couple the two volume estimates rather than discard their geometric relation.

For a sandwich built from a fixed arbitrary basis, taking $K=P$ makes the lower estimate for $|K|$ lose precisely $n!$, while the polar estimate is exact; taking $K=C$ reverses the roles. These are comparisons of the two volume bounds, not an assertion that this fixed basis is a maximal-determinant basis for each chosen body. For instance, a square can have a larger determinant using diagonal vertex vectors than using its coordinate half-edges. The comparisons expose the loss in the sandwich estimates; a sharp improvement needs additional geometric information.

\paragraph{6. An integral reformulation with all constants tracked.}
Let $\|x\|_K$ denote the gauge of $K$. Since the sublevel set $\{\|x\|_K\leq r\}$ is $rK$, layer-cake integration gives
\[
 \int_{\mathbb R^n}e^{-\|x\|_K}\,dx
 =n|K|\int_0^\infty e^{-r}r^{n-1}dr=n!|K|.
\]
Thus the conjecture is equivalent, for these particular gauge-exponential functions, to
\[
 \left(\int e^{-\|x\|_K}dx\right)
 \left(\int e^{-\|y\|_{K^\circ}}dy\right)\geq4^n n!.
\]
This is a restatement, not a proof from a general functional inequality. In particular the two functions are linked by polarity, not independent even log-concave functions. A theorem for another Legendre-dual normalization cannot be inserted without checking the transform and the factorials.

For the quadratic gauge potential $\varphi(x)=\|x\|_K^2/2$, its Legendre transform is $\varphi^*(y)=\|y\|_{K^\circ}^2/2$. To verify this, maximize $\langle x,y\rangle-\|x\|_K^2/2$ first over the radius $r=\|x\|_K$, using $\sup_{x\in K}\langle x,y\rangle=\|y\|_{K^\circ}$, and then maximize $r\|y\|_{K^\circ}-r^2/2$. The optimum is half the squared dual gauge. This supplies a legitimate dual pair, but still leaves the sharp integral lower bound to prove.

\paragraph{7. A smooth variational route and its limitation.}
For a smooth strictly convex body with positive support function $h$, polar coordinates in the polar body give
\[
 |K^\circ|=\frac1n\int_{S^{n-1}}h(u)^{-n}\,d\sigma(u).
\]
For an admissible support variation $h_t=h+tf$, differentiation yields
\[
 \left.\frac d{dt}\right|_0|K_t^\circ|
 =-\int_{S^{n-1}}f h^{-n-1}\,d\sigma.
\]
The usual first variation of volume is $\left.\frac d{dt}\right|_0|K_t|=\int f\,dS_K$, where $S_K$ is surface-area measure. Therefore a smooth stationary body for unrestricted symmetric support variations must satisfy, as even measures,
\[
 |K^\circ|\,dS_K=|K|\,h^{-n-1}d\sigma.
\]
This necessary equation can organize a local analysis. It does not establish a global minimum over all convex bodies, whose candidates include nonsmooth polytopes. For example, a Euclidean disk has area product $\pi^2$, whereas a square has product eight; stationarity at a smooth body cannot be taken as minimality. Compactness or existence of a minimizer alone also does not identify its value.

\paragraph{8. Diagnostics and outcome.}
The accompanying exact diagnostic checks beta-integral constants, normalized products for recursively constructed equality bodies in small dimensions, and the maximal-basis bound's factorial gap. It uses rational volumes and does not claim to exhaust all polytopes or certify a continuum of support variations.

\textbf{UNRESOLVED.} Exact product and convex-sum formulas prove the target for the recursively constructed equality family and propagate any known lower-dimensional subcases. The maximal-determinant argument proves a weaker universal bound. The first general gap is a coupled lower bound for $|K|$ and $|K^\circ|$ recovering the missing factor $n!$, or another argument controlling arbitrary centrally symmetric bodies beyond these decompositions.

\subsection{Sources}
\begin{itemize}
\item[S1]Karasev, R., 'Mahler's conjecture for some hyperplane sections', Israel J. Math. 241 (2021), 795-815.
\url{https://arxiv.org/pdf/1902.08971}.
\textit{Formulation source.}
\item[E1]H. Iriyeh and M. Shibata, Symmetric Mahler's conjecture for the volume product in the three dimensional case, arXiv:1706.01749v3, Theorem 1.1.
\url{https://arxiv.org/html/1706.01749}.
\textit{Primary theorem and dimension restriction checked 27 September 2026.}
\item[E2]The Symmetric Mahler Inequality in Dimension Three via Admissible Shadow Systems, arXiv:2605.13795, May 2026.
\url{https://arxiv.org/abs/2605.13795}.
\textit{Primary abstract inspected 27 September 2026: an alternative three-dimensional proof, not an all-dimensional resolution; full proof not independently audited.}
\end{itemize}
\end{document}
